\section{A bound for the entire global Weil form, with every prime retained}
\label{sec:global-weil-modulation}

\paragraph{Source and actual scope.}
The global input is the original trace formula of Alain Connes,
Caterina Consani and Matilde Marcolli,
\href{https://arxiv.org/abs/math/0703392v1}{\emph{The Weil proof and the
geometry of the adeles class space}}, author labels
\texttt{TraceformulaH1}, \texttt{fsharpinv}, \texttt{vanishV} and
\texttt{degmodifyV}. Its original real principal value is evaluated
in Alain Connes,
\href{https://arxiv.org/abs/math/9811068v1}{\emph{Trace formula in
noncommutative geometry and the zeros of the Riemann zeta function}},
Appendix II(30)--(37). The complete finite-place evaluation is WLC1--18;
the actual original theta source is ETC1--15. This calculation assembles
the entire global form and proves an all-frequency tail bound on it.
No finite list of test windows or numerical zero samples is used.
The cited global trace theorem remains the human theorem; the
identities and estimate below are derived with its full original terms.

\subsection{The original test space and the complete global form}

Let $\mathcal B_\infty$ be the smooth functions $h:(0,\infty)\to\mathbb C$
with every seminorm
\begin{equation}\tag{GWM1}
 p_{A,k}(h)=\sup_{x>0}(x^A+x^{-A})|(x\partial_x)^k h(x)|<\infty,
 \qquad A,k\in\mathbb Z_{\ge0}.
\end{equation}
This retains the original coordinate $x$ and both ends. For a character
$\chi$ of $U_N=(\mathbb Z/N\mathbb Z)^\times$, let $\chi_*$ be its
inducing primitive integer character, $c_\chi$ its conductor, and
$\chi_*(-1)=(-1)^{\epsilon_\chi}$, $\epsilon_\chi\in\{0,1\}$.
Use the actual adelic-class test $f(x,u)=h(x)\chi(u)$ with measure
$dx/x$ and unit-group Haar mass one. Its original sharp and correlation are
\begin{equation}\tag{GWM2}
 f^\sharp(x,u)=x^{-1}\overline{f(x^{-1},u^{-1})},\qquad
 (f\star f^\sharp)(y,u)=H_h(y)\chi(u),\qquad
 H_h(y)=\frac1y\int_0^\infty h(x)\overline{h(x/y)}\,dx.
\end{equation}
Indeed the compact integral is the integral of
$\chi(v)\chi(v^{-1}u)$, and the real factor in the convolution is
$(y/x)^{-1}dx/x=dx/y$. This proves GWM2 without changing measure.
Write
\begin{equation}\tag{GWM3}
 \mathcal M h(s)=\int_0^\infty h(x)x^s\frac{dx}{x},\quad
 M_h=\int_0^\infty|h(x)|^2dx=H_h(1),\quad
 H_h(y^{-1})=y\overline{H_h(y)}.
\end{equation}
The last identity follows by $x=yz$ in the conjugate integral.

The complete real-place term, with the original additive character
$e^{-2\pi ix}$, is
\begin{equation}\tag{GWM4}
 \begin{split}
 \mathcal A_\epsilon(H)={}&(\log(2\pi)+\gamma)H(1)\\
 &+\frac12\int_0^1\left[
 \frac{H(y)+y^{-1}H(y^{-1})-2H(1)}{1-y}
 +(-1)^\epsilon\frac{H(y)+y^{-1}H(y^{-1})}{1+y}
 \right]dy.
 \end{split}
\end{equation}
Here $\gamma$ is Euler's constant. To obtain this formula from Connes'
Appendix II(36), retain $d^\times u=du/(2|u|)$, substitute $y=|u|^{-1}$
separately on the two real signs, and retain the negative-sign character
$(-1)^\epsilon$. The positive part is
$\tfrac12\int H(y)|1-y|^{-1}dy$ with the subtraction
$\log\delta\,H(1)$ at $|1-y|=\delta$. Inverting the part $y>1$
gives $y^{-1}H(y^{-1})$ on $(0,1)$. Subtracting $2H(1)$ gives exactly
the displayed convergent integral, since
$\int_0^{1-\delta}(1-y)^{-1}dy=-\log\delta$.
The transformed cutoff differs from this symmetric cutoff by intervals
whose endpoint ratios tend to one, and hence by a term tending to zero.
The negative part has no singularity and yields the denominator $1+y$.
This retains the full original principal-value constant.

The complete global quadratic form is consequently
\begin{equation}\tag{GWM5}
 \begin{split}
 Q_\chi(h)={}&\delta_{\chi_*,1}\left[
 \mathcal Mh(0)\overline{\mathcal Mh(1)}+
 \mathcal Mh(1)\overline{\mathcal Mh(0)}\right]
 -\log|\mathfrak a|\,M_h\\
 &+\sum_{p\mid c_\chi}v_p(c_\chi)\log p\,M_h
 -\mathcal A_{\epsilon_\chi}(H_h)\\
 &-\sum_{p\nmid c_\chi}\log p\sum_{j\ge1}
 \left[\chi_*(p)^j H_h(p^j)
       +p^{-j}\chi_*(p)^{-j}H_h(p^{-j})\right].
 \end{split}
\end{equation}
The differential idele has $|\mathfrak a|=1$ over $\mathbb Q$ with these
original reference characters; its contribution has been retained in its
place. In particular at $\chi=1$ every rational prime remains present.
At an imprimitive modulus, primes dividing $N$ but not $c_\chi$ still
occur; the original comparison is
$L(s,\chi)=L(s,\chi_*)\prod_{p\mid N,\ p\nmid c_\chi}
(1-\chi_*(p)p^{-s})$. No zero value of the imprimitive character replaces
a primitive value in GWM5. The conductor sum is precisely the negative
of the ramified local distribution, not an additional Euler factor.

Every term is defined on GWM1, including the infinite sum over all primes.
For $n\ge1$ and integer $A\ge2$, put
$C_A=\sup_x\max(x^A,x^{-A})|h(x)|$. Splitting the integral in GWM2
at $1$ and $n$ gives the explicit bound
\begin{equation}\tag{GWM6}
 |H_h(n)|\le C_A^2 n^{-A}
 \left[\frac{n^{-1}}{2A+1}+1-n^{-1}+\frac1{2A-1}\right]
 \le2C_A^2 n^{-A}.
\end{equation}
The three terms are the integrals on $(0,1)$, $(1,n)$, and $(n,\infty)$,
respectively. Thus all prime powers, with any additional fixed power of
$\log n$, converge absolutely by comparison with
$\sum_{n\ge2}(\log n)^k n^{-A}$. The estimates after multiplicative
differentiation prove the same statements for the smooth topology GWM1.
The real integral converges at zero by the same endpoint bounds and at
one by subtraction and smoothness. The endpoints converge separately.

\subsection{Exact spectral evaluation of the retained real contribution}

The full Gamma contribution in GWM4 can also be evaluated, without
replacing the original arithmetic function:
\begin{equation}\tag{GWM7}
 -\mathcal A_\epsilon(H_h)=\frac1{2\pi}\int_{\mathbb R}
 \left[\Re\psi\left(\frac{1/2+it+\epsilon}{2}\right)-\log\pi\right]
 |\mathcal Mh(1/2+it)|^2\,dt.
\end{equation}
Here $\psi=\Gamma'/\Gamma$. The multiplier is an evaluated summand
of GWM5, not a substitute function for $\zeta$ or $L(s,\chi_*)$.
To prove it, set $u=\log x$ only inside the integration argument.
Fourier inversion of the correlation in GWM2 gives
\begin{equation}\tag{GWM8}
 H_h(y)=\frac1{2\pi}\int_{\mathbb R}
 |\mathcal Mh(1/2+it)|^2 y^{-1/2-it}\,dt,\qquad
 M_h=\frac1{2\pi}\int_{\mathbb R}|\mathcal Mh(1/2+it)|^2dt.
\end{equation}
Indeed the function whose Fourier integral occurs is
$e^{u/2}h(e^u)$, and its squared norm is exactly
$\int |h(x)|^2dx$; this records both the factor and the Jacobian.
It is Schwartz by GWM1. Substitution into GWM4 gives, for each $t$,
\[
 \log(2\pi)+\gamma+
 \int_0^1\left[
 \frac{y^{-1/2}\cos(t\log y)-1}{1-y}
 +(-1)^\epsilon\frac{y^{-1/2}\cos(t\log y)}{1+y}\right]dy.
\]
With $z=y^2$ the integral is
\[
 \int_0^1\frac{z^{\alpha-1}\cos((t/2)\log z)-1}{1-z}\,dz
 +\frac12\int_0^1\frac{1-z^{-1/2}}{1-z}\,dz,
 \qquad\alpha=\frac{1/2+\epsilon}{2}.
\]
The digamma series gives the first integral
$-\Re\psi(\alpha+it/2)-\gamma$ and the second is $-\log2$.
For the latter, put $z=r^2$ to get $-\int_0^1(1+r)^{-1}dr$.
Thus the full expression is $\log\pi-\Re\psi(\alpha+it/2)$,
proving GWM7 with all constants accounted for. Interchanging integrals
is justified at $y=1$ by a bound $C(1+t^2)$ after subtraction, at $0$
by the integrable factor $y^{-1/2}$, and in $t$ by Schwartz decay.

\subsection{A proved lower bound on the entire form for all large frequencies}

Keep the original function and introduce its explicitly invertible
modulation $h_T(x)=x^{iT}h(x)$, $T\in\mathbb R$. The exact relations are
\begin{equation}\tag{GWM9}
 H_{h_T}(y)=y^{iT}H_h(y),\qquad
 \mathcal Mh_T(s)=\mathcal Mh(s+iT),\qquad M_{h_T}=M_h.
\end{equation}
Define the following quantities from the full original $h$:
\begin{equation}\tag{GWM10}
 \begin{gathered}
 K_h=\int_0^\infty|xh'(x)+\tfrac12h(x)|^2dx,\qquad
 A_{0,h}=\int_0^\infty|h(x)|\frac{dx}{x},\qquad
 A_{1,h}=\int_0^\infty|h(x)|dx,\\
 P_{\chi,h}=2\sum_{p\nmid c_\chi}\log p\sum_{j\ge1}|H_h(p^j)|,
 \qquad\alpha_\chi=\frac{1/2+\epsilon_\chi}{2}.
 \end{gathered}
\end{equation}
The sum in $P_{\chi,h}$ is over every prime power, including when $h$
has noncompact support. Its convergence was proved in GWM6. For every
$|T|\ge4\alpha_\chi$, the complete global form satisfies
\begin{equation}\tag{GWM11}
 \begin{split}
 Q_\chi(h_T)\ge{}&M_h\left[
 \log\frac{|T|}{4}-\frac1{\alpha_\chi}-\log\pi
 +\sum_{p\mid c_\chi}v_p(c_\chi)\log p-\log|\mathfrak a|\right]\\
 &-\frac{4K_h}{T^2}\log\frac{|T|}{4\alpha_\chi}
 -P_{\chi,h}-2\delta_{\chi_*,1}A_{0,h}A_{1,h}.
 \end{split}
\end{equation}
This is a bound on all of GWM5, not merely its archimedean part.

Here is a proof of every estimate used. For $\alpha>0$ and real $b$,
the digamma series, with the harmonic-number constant retained, gives
\[
 \Re\psi(\alpha+ib)=
 \lim_{m\to\infty}\left[\log m-
 \sum_{k=0}^{m-1}\frac{k+\alpha}{(k+\alpha)^2+b^2}\right].
\]
For $r(x)=(x+\alpha)/((x+\alpha)^2+b^2)$,
$|\sum_{k=0}^{m-1}r(k)-\int_0^m r(x)dx|
\le\int_0^m|r'(x)|dx$ by integration on every unit interval.
If $|b|\le\alpha$, this total variation is $r(0)\le1/\alpha$.
If $|b|>\alpha$, it is
$1/|b|-\alpha/(\alpha^2+b^2)\le1/\alpha$.
The integral is
$\tfrac12\log(((m+\alpha)^2+b^2)/(\alpha^2+b^2))$.
Consequently the unconditional pointwise bound is
\begin{equation}\tag{GWM12}
 \Re\psi(\alpha+ib)\ge
 \frac12\log(\alpha^2+b^2)-\frac1\alpha.
\end{equation}
In GWM7 for $h_T$, put $u=t+T$ without altering the original $dt/(2\pi)$
factor. On $|u|\le|T|/2$ one has $|t|\ge|T|/2$, so GWM12 is at least
$\log(|T|/4)-1/\alpha_\chi$. Everywhere it is at least
$\log\alpha_\chi-1/\alpha_\chi$. The exact Parseval identities give
\begin{equation}\tag{GWM13}
 \frac1{2\pi}\int_{\mathbb R}u^2|\mathcal Mh(1/2+iu)|^2du=K_h,
 \qquad
 \frac1{2\pi}\int_{|u|>|T|/2}|\mathcal Mh(1/2+iu)|^2du
 \le\frac{4K_h}{T^2}.
\end{equation}
For the first identity, differentiating $e^{u/2}h(e^u)$ gives exactly
$e^{u/2}(xh'(x)+h(x)/2)$; its norm has measure $dx$, as in GWM8.
Subtracting the two lower bounds just obtained gives the explicit
tail defect $4K_h T^{-2}\log(|T|/(4\alpha_\chi))$.
Each paired prime summand in GWM5 is
$2\Re(\chi_*(p)^j p^{ijT}H_h(p^j))$, by GWM3 and GWM9,
so its total possible negative contribution is bounded by $P_{\chi,h}$.
The two original endpoint integrals contribute at least
$-2\delta_{\chi_*,1}A_{0,h}A_{1,h}$. The conductor and differential
terms are unchanged by GWM9. Combining these estimates proves GWM11.

There is also an explicit finite threshold, with no numerical certificate.
For nonzero $h$, let
\begin{equation}\tag{GWM14}
 \begin{split}
 T_{\chi,h}=\max\Biggl\{4\alpha_\chi,\quad4\exp\Biggl[
 1+\frac1{\alpha_\chi}+\log\pi
 -\sum_{p\mid c_\chi}v_p(c_\chi)\log p+\log|\mathfrak a|\\
 {}+\frac{P_{\chi,h}+2\delta_{\chi_*,1}A_{0,h}A_{1,h}}{M_h}
 +\frac{K_h}{8\alpha_\chi^2\exp(1)M_h}
 \Biggr]\Biggr\}.
 \end{split}
\end{equation}
For every $|T|\ge T_{\chi,h}$ one has $Q_\chi(h_T)\ge M_h>0$.
Indeed the maximum of $\log(t/(4\alpha))/t^2$ on $t\ge4\alpha$
is $1/(32\alpha^2\exp(1))$, as its derivative has its sole maximum
at $t=4\alpha\sqrt{\exp(1)}$. Substitution into GWM11 proves the
assertion. The threshold depends on the entire original amplitude;
no uniform threshold over all $h$ has been asserted.

\subsection{Global character tuples and the original support carrier}

For $f(x,u)=\sum_{\chi\in\widehat U_N}h_\chi(x)\chi(u)$, direct compact
character orthogonality in the convolution gives
\begin{equation}\tag{GWM15}
 Q(f)=\sum_{\chi\in\widehat U_N}Q_\chi(h_\chi).
\end{equation}
The mixed $\chi,\psi$ term has compact integral
$\int\chi(v)\overline{\psi(v)}du=\delta_{\chi,\psi}$; the sharp
preserves the character label. Thus GWM15 is an exact global
diagonalization and leaves every ramified character and its conductor.
For a nonzero finite tuple, taking the maximum of GWM14 over its
nonzero entries proves
$Q(|\cdot|^{iT}f)\ge\sum_\chi M_{h_\chi}>0$ for every $|T|$ beyond
that threshold. All primes remain in each of its global constants.

Let $L$ be the original nontrivial bounded distributive lattice and
$\mathcal B$ the original finite-character amplitude algebra with
the functions GWM1. Its full carrier and maps are
\begin{equation}\tag{GWM16}
 \begin{gathered}
 G_L(\mathcal B)=\{(0,\lambda):\lambda\in L\}
                      \cup(\mathcal B\times\{1_L\}),\\
 \Delta_{iT}^{\uparrow}(f,\lambda)=(|\cdot|^{iT}f,\lambda),\qquad
 Q^{\uparrow}((f,\lambda),(g,\mu))=(Q(f,g),\lambda\wedge\mu).
 \end{gathered}
\end{equation}
The first map is an invertible complex-linear and convolution
algebra map: $|xy|^{iT}=|x|^{iT}|y|^{iT}$ proves the convolution
identity under its integral. It commutes with sharp since
$\overline{|x^{-1}|^{iT}}=|x|^{iT}$. Its inverse is $\Delta_{-iT}$.
It keeps each label, including $\tau$ and $e$. The second map is
well-defined because nonzero amplitude requires both input labels
to be top. Sesquilinearity and finite distributivity of the lattice
prove additivity in each supported argument and the original scalar
rule. A vanished top-labelled global pairing has value $e$, not
$\tau$. At every support prime $J$, each label map restricts to
the same meet in $L_J$. No measure or positivity ordering on an
arbitrary lattice is added.

These are maps on the original test algebra. No descent of
$\Delta_{iT}$ to the arithmetic spectral quotient is assumed:
the actual adelic image need not be preserved by modulation.
The global kernel built from ETC's theta source computes that
receiving defect and reconstructs all compact tests. In particular,
the large-frequency result is a proved statement about the full
form on specified modulated tests, not a conclusion about its sign
on every test or about RH.

\subsection{One global bound for the entire original adelic theta family}

Restore both original parameters from ETC, before taking their image:
\begin{equation}\tag{GWM17}
 \begin{gathered}
 k\in\bigoplus_p\mathbb Z,\qquad a>0,\qquad
 N(k)=\prod_p p^{k_p},\qquad b=\frac a{N(k)},\\
 \xi_{k,a}(x)=\eta(x_\infty/a)\prod_p1_{p^{k_p}\mathbb Z_p}(x_p),\qquad
 \eta(x)=\pi x^2(\pi x^2-3/2)e^{-\pi x^2},\\
 \mathcal E\xi_{k,a}(y,u)=f_b(y)
      =2\sum_{n\ge1}\eta(ny/b)=f_1(y/b).
 \end{gathered}
\end{equation}
ETC3--9 proves the two original additive moments vanish, the complete
Fourier and rational restriction formulas, and the rapid bounds at both
ends. Thus every $f_b$ belongs to GWM1 and to the actual image
$\mathcal V$ of the adelic source. Their positivity is also exact:
for $y\ge1$ every summand of $f_1(y)$ is positive because
$\pi n^2y^2>3/2$; for $0<y<1$, the proved Poisson formula
$f_1(y)=y^{-1}f_1(y^{-1})$ gives positivity. Scaling proves $f_b>0$.
No theta function is inferred from the whole absolute base by this
argument; GWM17 is the specified original tensor and its proved map.

The full original Mellin formula is
\begin{equation}\tag{GWM18}
 \mathcal M f_b(s)=2\left(\frac a{N(k)}\right)^s\zeta(s)\left[
 \frac{\pi^2}{2\pi^{(s+4)/2}}\Gamma\left(\frac{s+4}{2}\right)
 -\frac{3\pi}{4\pi^{(s+2)/2}}\Gamma\left(\frac{s+2}{2}\right)\right],\qquad
 \mathcal M f_b(0)=\frac14,\quad\mathcal M f_b(1)=\frac b4.
\end{equation}
The expression on the right agrees first for $\Re s>1$ by absolute
summation and Gaussian integration, and then by continuation with the
entire Mellin integral on the left, as proved in ETC10--15. Both
Gaussian summands, the original $\zeta$, all Gamma factors, the power
of $\pi$, and the real dilation factor are retained. In particular
it vanishes at every nontrivial zero of the original $\zeta$ with
at least its original multiplicity. Alternatively the source ideal
calculation \texttt{vanishV} directly gives $Q_1(f_b)=0$.
Using every term of GWM5 and these two endpoints gives the global identity
\begin{equation}\tag{GWM19}
 \mathcal A_0(H_{f_b})+
 2\sum_p\log p\sum_{j\ge1}H_{f_b}(p^j)
 +\log|\mathfrak a|\,M_{f_b}
 =\frac b8,\qquad |\mathfrak a|=1.
\end{equation}
Here the correlation is real and positive and the prime sum includes
every prime power. Its absolute convergence was proved in GWM6.
This is an exact global balance, not a bound obtained from a cutoff.

Changing variables in each original integral proves all scaling factors:
\begin{equation}\tag{GWM20}
 \begin{gathered}
 H_{f_b}(y)=bH_{f_1}(y),\quad M_{f_b}=bM_{f_1},\quad
 K_{f_b}=bK_{f_1},\quad P_{1,f_b}=bP_{1,f_1},\\
 A_{0,f_b}=A_{0,f_1},\qquad A_{1,f_b}=bA_{1,f_1},\qquad
 Q_1(x^{iT}f_b(x))=bQ_1(x^{iT}f_1(x)).
 \end{gathered}
\end{equation}
For the last equality, GWM9 proves the correlation scaling, while the
two Mellin endpoints scale by $b^{iT}$ and $b^{1+iT}$; their conjugate
product retains exactly the factor $b$. Every other summand of the
complete global form scales by $b$. The threshold GWM14 is therefore
exactly the same for every $a>0$ and every finite divisor $k$:
\begin{equation}\tag{GWM21}
 T_{1,f_b}=T_{1,f_1},\qquad
 Q_1(x^{iT}f_b(x))\ge bM_{f_1}>0
 \quad\text{for every }|T|\ge T_{1,f_1}.
\end{equation}
This proves a global result simultaneously for the whole original
source family and the whole frequency tail, with the factor $a/N(k)$
kept. It does not recover $a,k$ separately from a scalar map that ETC6
proved depends on their ratio.

There is an exact geometric consequence at the arithmetic receiving
quotient. Let $\mathcal W$ be the closure of $\mathcal V$ in the original
strong test topology, and $q:\mathcal B\to\mathcal B/\mathcal W$ its
linear quotient. The bounds above show that the full pairing is
continuous. The cited source ideal calculation makes $\mathcal V$
orthogonal to every test; continuity gives the same for $\mathcal W$.
Consequently
\begin{equation}\tag{GWM22}
 \mathfrak D_T:\mathcal V\longrightarrow\mathcal B/\mathcal W,
 \qquad v\longmapsto q(|\cdot|^{iT}v),\qquad
 \ker\mathfrak D_T=\mathcal V\cap\Delta_{-iT}\mathcal W.
\end{equation}
This is an actual linear map with the displayed full kernel, not an
assumed descended modulation. It has $\mathfrak D_0(f_b)=0$, but
$\mathfrak D_T(f_b)\ne0$ for every $|T|\ge T_{1,f_1}$, by GWM21:
a zero quotient class would have zero pairing. Hence modulation
does not preserve this original arithmetic quotient on that entire
frequency tail. Its exact full-label lift is
$(v,\lambda)\mapsto(\mathfrak D_Tv,\lambda)$; the source
$(f_b,1_L)$ has image $e$ at $T=0$ and a nonzero top-labelled
amplitude for every frequency in GWM21. No supported zero is
identified with $\tau$, and no omitted local stratum is used to
deduce this global non-descent result.
